% Replace the sample sections with your own policies and procedures.
\section{Purpose and scope}
This handbook describes how the organization protects its systems and data.
It applies to every employee, contractor and service provider who has access
to company systems.

\begin{policybox}
  All company data must be classified, stored in an approved system and
  shared only with people who need it for their work.
\end{policybox}

\subsection{Roles and responsibilities}
Every role has a clear owner. The table lists who decides, who operates and
who must be told.

\begin{table}[H]
  \centering
  \begin{tabularx}{\linewidth}{@{}lXl@{}}
    \toprule
    \textbf{Role} & \textbf{Responsibility} & \textbf{Contact} \\
    \midrule
    Security lead & Owns this handbook and approves exceptions. & Security team \\
    System owner & Keeps access lists and patches up to date. & Platform team \\
    Employee & Follows the policies and reports incidents. & Everyone \\
    \bottomrule
  \end{tabularx}
\end{table}

\section{Access management}
Access is granted by role and reviewed every quarter. Shared accounts are not
allowed; every action must be traceable to one person.

\begin{procedurebox}
  \begin{steps}
    \item The manager files an access request naming the system and the role.
    \item The system owner checks that the role matches the job.
    \item The platform team grants the access and records the ticket number.
    \item The employee confirms access within two working days.
  \end{steps}
\end{procedurebox}

\begin{warningbox}
  Never share credentials in chat or email. If you suspect a password is
  exposed, change it first and report it afterwards.
\end{warningbox}

\section{Incident response playbook}
When something looks wrong, the first goal is to contain the problem, and the
second is to understand it. The flow below shows the path of every incident.

\begin{mermaid}[style=editorial,width=0.9\linewidth,
    caption={Incident response flow},pos=H]
  flowchart LR
  A[Detect] --> B[Triage] --> C[Contain] --> D[Recover] --> E[Review]
\end{mermaid}

The on-call engineer can list recent failed sign-ins with the command below.\par\vspace{2pt}

\begin{code}[lang=bash]
grep "authentication failure" /var/log/auth.log | tail -n 20
\end{code}

\begin{notebox}
  Record every action and its time in the incident ticket while you work. The
  review depends on that record.
\end{notebox}

\section{Decision records}
Significant choices are recorded so that later readers know why they were made.

\begin{decision}{DR-004}{Adopt hardware security keys}
  \decisionpart{\lblContext}
  Phishing attempts rose during the last two quarters and one-time codes can
  still be intercepted.

  \decisionpart{\lblDecision}
  All staff with access to production systems will use hardware security keys
  for sign-in from the first of January.

  \decisionpart{\lblConsequences}
  Each person needs two keys and a spare is kept in the office safe. Support
  must handle key loss within one working day.
\end{decision}

\clearpage
\section*{\lblGlossary}
\addcontentsline{toc}{section}{\lblGlossary}
\begin{termlist}
  \item[Classification] The label that tells readers how widely a document may be shared.
  \item[Incident] Any event that harms or could harm the confidentiality, integrity or availability of company data.
  \item[Runbook] A short, tested list of steps for one recurring task.
  \item[Least privilege] Giving a person only the access they need for their current work.
  \item[Review] A scheduled check that a policy, a role or an access list is still correct.
  \item[System owner] The person accountable for one system and its access list.
\end{termlist}
